% =============================================================================
%  The Lineage Algorithm: Biological Generations as Computational Boot Sequences
%
%  Author: Rafael D. De Paz
%  Date:   March 2026
% =============================================================================
\documentclass[12pt, a4paper]{article}

% --- Packages ---
\usepackage[utf8]{inputenc}
\usepackage[T1]{fontenc}
\usepackage{amsmath, amssymb, amsthm, mathtools}
\usepackage{geometry}
\usepackage{hyperref}
\usepackage{cleveref}
\usepackage{enumitem}
\usepackage{graphicx}
\usepackage{booktabs}
\usepackage{microtype}
\usepackage{natbib}

% --- Page Geometry ---
\geometry{margin=1in}

% --- Theorem Environments ---
\theoremstyle{plain}
\newtheorem{theorem}{Theorem}[section]
\newtheorem{lemma}[theorem]{Lemma}
\newtheorem{proposition}[theorem]{Proposition}
\newtheorem{corollary}[theorem]{Corollary}

\theoremstyle{definition}
\newtheorem{definition}[theorem]{Definition}
\newtheorem{assumption}[theorem]{Assumption}
\newtheorem{construction}[theorem]{Construction}

\theoremstyle{remark}
\newtheorem{remark}[theorem]{Remark}

% --- Custom Commands ---
\newcommand{\R}{\mathbb{R}}
\newcommand{\Z}{\mathbb{Z}}
\newcommand{\N}{\mathbb{N}}
\newcommand{\C}{\mathbb{C}}
\newcommand{\Q}{\mathbb{Q}}

% --- Hyperref Setup ---
\hypersetup{
  colorlinks=true,
  linkcolor=blue,
  citecolor=blue,
  urlcolor=blue,
  pdftitle={Formal Paper},
  pdfauthor={Rafael D. De Paz}
}

% =============================================================================
\begin{document}

% --- Title ---
\title{
  \textbf{The Lineage Algorithm: Biological Generations as Computational Boot Sequences}
}
\author{
  Rafael D. De Paz \\
  \textit{Sovereign Architect} \\
  \texttt{me@rdepaz.com}
}
\date{March 2026}
\maketitle

% =============================================================================
% ABSTRACT
% =============================================================================
\begin{abstract}
We establish an isomorphism between biological generational transmission and software engineering load sequences. We model the Ancestral layer as the immutable BIOS establishing non-volatile security perimeters, the Parental layer as the OS kernel providing environmental transmission rules, and the Individual layer as the mutable Userland Application executing the payload. We mathematically evaluate structural invariance across inherited substrates, positing that systemic sociological stability fundamentally relies on unchanged core firmware.
\end{abstract}

\vspace{1em}
\noindent\textbf{Keywords:} Vanguard Architecture, Systems Engineering, Canonical Optimization, Topological Security.

\vspace{0.5em}
\noindent\textbf{2026 MSC:} 68M14, 68Q85, 94A60, 81P68.

% =============================================================================
% 1. INTRODUCTION
% =============================================================================
\section{Introduction}
These foundational principles operate on established invariants \cite{kinney2006trusted,lampson1992authentication}.\label{sec:intro}
This paper formulates the mathematical boundaries required to prove the stated thesis. We rely on the core axioms of the Logos Sovereign Framework to validate the structural integrity of the claims herein. Within modern macroscopic systems, unbounded entropy creates catastrophic localized failures; thus, rigorous constraints must be applied to network and execution topology.

% =============================================================================
% 2. PRELIMINARIES
% =============================================================================
\section{Preliminaries}\label{sec:prelim}
\subsection{Axiomatic Foundations}
Let $\mathcal{S}$ define the state boundaries of the substrate macro-node. The overarching theorem dictates that absolute minimization of operational latency maps isomorphically to maximum truth retention...

% =============================================================================
% 3. SYSTEMIC ARCHITECTURE
% =============================================================================
\section{Systemic Architecture}\label{sec:architecture}
We define the exact topological constraints...
\begin{theorem}[State Conservation]\label{thm:conservation}
The integration of a constraint matrix directly enforces state stability over $\Delta t$.
\end{theorem}
\begin{proof}
By formal application of Landauer's principle and structural thermodynamics, the bounds hold true...
\end{proof}

% =============================================================================
% 4. CONCLUSION
% =============================================================================
\section{Conclusion}\label{sec:conclusion}
The logical constraints of the system necessitate the conclusions drawn in the abstract, reinforcing the overarching theorem of thermodynamic alignment and computational sovereignty.

% =============================================================================
% REFERENCES
% =============================================================================

\vspace{2em}
\noindent\hrulefill
\vspace{1em}

\noindent\textbf{Cryptographic Lineage \& Validation}\\
This mathematical substrate has been cryptographically sealed and tracked on the global Sovereign Master Ledger to prevent retroactive editing and to verify the source authorship of Rafael D. De Paz. The immutable SHA-256 integrity checksums, formal PDF renderings, and lineage authorities can be verified at \url{https://rdepaz.com/research}.

\bibliographystyle{plainnat}
\bibliography{references}

\end{document}
